\documentclass[notumble,combine]{leaflet}
\usepackage[dvipsnames,usenames]{color}
\usepackage{graphicx}
\usepackage{alltt}
\usepackage{url}
\usepackage{ascmac}
\usepackage{comment}
\usepackage{here}
\usepackage{wrapfig}
\usepackage{floatflt}

\makeatletter
% Section color
%% from leaflet.cls
\renewcommand\section{\@startsection{section}{1}{\z@}%
  {-3.5ex \@plus -.75ex}%
  {1ex} %{1.5ex}%
  {\normalfont\large\sectfont\color{NavyBlue}}}
\renewcommand\subsection{\@startsection{subsection}{2}{\z@}%
  {-2.5ex plus -.5ex}%
  {1\p@} %{1ex}%
  {\normalfont\normalsize\sectfont\color{Green}}}
\makeatother

\graphicspath{{figures/}} 

\title{
	\includegraphics[width=\textwidth]{kbug-logo}\\
        \parbox[c]{2cm}{\resizebox{2cm}{!}{\textcolor{red}{\bf{@}}}}\\
	\includegraphics[width=6cm]{osc2013-kyoto-logo}\\
	\vfill
	\parbox[c]{\textwidth}{\includegraphics[width=\textwidth]{MyKyotoIchiba.eps}}\\
        \vfill
	\resizebox{\linewidth}{!}{\bf 関西*BSDユーザ会}\\
        \resizebox{\linewidth}{!}{{\bf \url{http://www.kbug.gr.jp/}}}\\
        \includegraphics[width=3cm]{kbug.gr.jp.eps}\\
        \parbox[c]{1cm}{\resizebox{1cm}{!}{\textcolor{Orange}{\bf{@}}}}
        \resizebox{\linewidth}{!}{2013年8月2日(金),3日(土)}\\
}



\date{}

\begin{document}
\maketitle
\thispagestyle{empty}
\pagebreak{}
\section{関西*BSDユーザ会(K*BUG)ってなぁに?}
関西 *BSD ユーザ会 (Kansai *BSD Users Group; K*BUG)とは、
BSD系OSのユーザ同士の情報交換のための{\Large\em \textcolor{red}{場}}で、1999年に設立されました。

年に数度の勉強会やイベントの参加、そして{\Large\em \textcolor{red}{飲み会}}を行っています。

\subsection{K*BSDの基本理念}
K*BUGの基本理念はのとおりです。

\fbox{\begin{minipage}{\textwidth}
\begin{itemize}
\item 場の提供を目的とする
\item 人のケツは叩くが足は引っ張らない
\item 来るものは拒まず、猿ものは追わず
\item だれでも役員になれる。誰でも役員は止められる
\end{itemize}
\end{minipage}}

少し難しく感じるかもしれないですが、
\textcolor{red}{\bf\Large BSDへの愛と情熱}
があれば、
あなたが
\textcolor{blue}{やりたいと思うことができる場}
がK*BUGなのです。

\section{K*BUG Keywords}
K*BUGでは、以下のようなKeywordに関する発表や展示を行いました。

{\large\bf\textcolor{red}{{BSD}}} (FreeBSD (PC-BSD), NetBSD, OpenBSD, DragonlyBSD, {\footnotesize Mac OS X(?), iOS(??) }, ...
の更新情報 / 利用方法 / ...,
arch {\footnotesize (i386, amd64, macppc, landisk, zaurus, wzero3, hpcmips, fonera, netwalker, ipaq, fonera, ...)},
kernel hackや新技術 (DTrace, ZFS, シリアルドライバ, ...),
FreeBSD portsってなに?,
Security, 
Prolog, Lisp,
Squeak, Scratch,
OpenCV,
ContaoCMS, 
hardware (UPS, HDD, Server, ...),
ものづくり,
電子工作,
ニコ動技術部 (猫耳サーボ % (USB audio servo controler)
, 鼻ホタル, ...),
雑誌付録基板 ,
Physical Computing (Gainer (gainerm-lib), Arduino, ...),
木彫, でーもんくん, でもんむし君 , \\
\hfill{\textcolor{red}{{\bf\large 飲み会}}}, and {\em\textcolor{red}{{ more with }}{\Large\bf\textcolor{red}{{YOU!!}}}}

\newpage
\section{K*BUGの活動}
最近の活動はこんな感じです。
もし、興味のあるテーマがあったら、
\begin{center}
{\huge\bf \textcolor{red}{遊びに来てね!!}}
\end{center}

\subsection{2013年06月22日(土)@大阪 第3回研究会}
\begin{itemize}
\item 会場ナレッジキャピタルに関して
\item FreeBSD ports 与太話
\item Raspberry Pi の話
\item ARMのおはなし 
\end{itemize}

\subsection{2013年4月20日(土)@大阪 第2回研究会}

\subsection{2013年4月6日(土)\\\hfill  第13回イベント@奈良高専(第1.5回研究会)}
\begin{itemize}
\item 付録基板で遊ぼ (BSD編)
\item Open Source Software (OSS) の魅力 (パネル)
\end{itemize}

\subsection{2013年2月16日(土)@大阪 第1回研究会}
\begin{itemize}
\item カーネル/VM勉強会@関西4に合流
\end{itemize}

\subsection{2012年12月8日(土)@京都 \\\hfill 定期総会+第7回研究会+忘年会}
\begin{itemize}
\item 関西FPGA・DE0勉強会に行ってきました
\item PCIブリッジのコンフィグレーション
\item FreeBSD portsと暮らす
\item Contao Open Source CMSでコンテンツ管理
\end{itemize}

\subsection{2012年11月9日(金), 10日(土) @大阪 \\\hfill KOF2012}
\begin{itemize}
\item BSDなひととき(JNUGセミナーに合流)
  \begin{itemize}
  \item K*BUGってなあに?@KOF2012
  \end{itemize}
\item 魅惑の付録基板の世界
\end{itemize}

  \begin{center}
    \includegraphics[width=4cm]{KOF2012Booth}
  \end{center}

\subsection{2012年10月20日(土)@大阪 第6回研究会}
\begin{itemize}
\item NetBSD 6.0Rの変更点
\end{itemize}

\subsection{2012年 8月18日(土)@京都 第5回研究会}
\begin{itemize}
\item 新しいレノボPC X230買いました
\item NetBSD on kobo
\item FREQUPS UPF F and NetBSD
\item Multi Thread Tiny BASIC on PIC32\\
  and RetroBSD
\item androidのアプリ
\item Buffaloルータのシリアル
\end{itemize}

\subsection{2012年 8月3日(金), 4日(土) @京都 \\\hfill OSC2012 Kansai@Kyoto}

\begin{itemize}
\item OpenBSDでNAT64
\item RetroBSD展示
\item でーもんくんの行進:
  miniでーもんくん, 
  ダンボールでーもんくん % Arduino
  よろこび棒 OSC2012, ... % OSC2012 Kansai@Kyoto版
\item K*BUG POV (失敗) %with Japanino
\item チラシとでもんむし君シール配布
\end{itemize}

\newpage
\section{これからの関連イベント(予定)}
K*BUGでは、2ヶ月に1度程度の頻度で勉強会を行っています。
また、主に関西のイベントで、展示などを行っています。

現在、予定されているイベントは、以下のとおりです。
詳細に関しては、K*BUGのWebページをご確認ください。

\begin{itemize} 
\item 
\textcolor{red}{\Large\em もうすぐ!!}
  2013年8月24日(土)@グランフロント大阪ナレッジキャピタル7F ナレッジサロン\\
  第4回研究会
\item 2013年10月19日(土)@AXE会議室(京都)\\
  第5回研究会
\item \textcolor{Gray}{2013年11月8日(金),9日(土) KOF2013} (参加未定)
\item 2013年12月7日(土)@大阪学院大学2号館02-03-02教室\\
  第15回定期総会 + 第6回研究会
\end{itemize}

\vfill
\begin{minipage}{\textwidth}
\begin{boxnote}

\section{kbug-usersメーリングリスト}
K*BUGでは、K*BUGメンバーの情報交換や、イベントなどの情報伝達用に
kbug-usersメーリングリストを用意しています。

  \begin{floatingfigure}[r]{2.5cm}{
      \includegraphics[width=2.5cm]{mailto.kbug-users.eps}
  }
  \end{floatingfigure}

K*BUGメンバーは、基本的にはこのメーリングリストを読んでいることが期待
されます。

購読は、右のQRコードを使うか、このURL
\footnote{\url{http://www.kbug.gr.jp/maillist.html}}
をご参照ください。

\end{boxnote}
\end{minipage}

\newpage
\section{K*BUG at OSC2013 Kansai@Kyoto}
  \begin{floatingfigure}[r]{2cm}{
      \includegraphics[width=2cm]{qml.610t.org.freebsd.OSC2013Kyoto.eps}
  }
  \end{floatingfigure}

K*BUGは、今回OSC2013 Kansai@Kyotoに以下のような活動で参加しています。

右のQRコードで、K*BUG OSC2013 Kansai@Kyoto特設ページ
\footnote{\url{http://qml.610t.org/FreeBSD/OSC2013Kyoto.html}}
にアクセスできます。

\subsection{K*BUGブース}
BSDを使って、以下の様な展示をしています。

\begin{itemize}
\item こんなにあるよ!! BSD色々展示  
\item 真夏のレシピ: 基板の気分
\item 
\textcolor{red}{\Large\em 配布!!}
でもんむし君シールとチラシ%\footenote{\url{}} 
\end{itemize}

\subsection{関連セミナー}
BSDに関して、みんなでお話しましょう!!
来場者によるライトニングトークなどの話題提供を歓迎します。
発表希望者は、お気軽にひとことどうぞ。

\subsubsection{NetBSDユーザのご紹介} % by 日本NetBSDユーザーグループ}
\begin{itemize}
\item 日時と会場: 2013/8/2(金) 15:15-16:00 at 会議室C
\item 発表: 蛯原 純 (The NetBSD Project)
\item 内容:
  \begin{itemize}
  \item NetBSD/luna68kはmlterm-fbでTwitterの夢を見るか(仮)
  \item NetBSD だと USB 3.0 がここまで動いてます(キリッ
    \hfill など
  \end{itemize}
\end{itemize}

\subsubsection{関西*BSDユーザ会のご紹介}
\begin{itemize}
\item 日時と会場: 2013/8/2(金) 16:15-17:00 at 会議室C
\item 発表: 関西*BSDユーザ会有志一同
\item 内容:
  \begin{itemize}
  \item K*BUGってなぁに?
  \item 真夏のレシピ: 基板の気分
    \hfill など
  \end{itemize}
\end{itemize}

\vfill
\verb+$Date: 2013/07/27 07:40:14 $ $Revision: 1.25 $+

\end{document}
